% from ega3-1-fr.tex
\oldpage[0\textsubscript{III}]{50}
à la restriction à la catégorie $C(X)$ du foncteur cohomologique défini de
même sur la catégorie des \emph{faisceaux de groupes abéliens} sur $X$
(G, II, 7.2.1).
\end{env}

\begin{env}[12.1.2]
\label{0.12.1.2-fr}
Posons $A=\Gamma(X,\mathcal O_X)$. Comme tout élément de $A$ définit un
endomorphisme du groupe abélien $\Gamma(X,\mathcal F)$, il définit par
fonctorialité un endomorphisme de $\delta$-foncteur de
$H^\bullet(X,\mathcal F)$; ces endomorphismes définissent sur chacun des
$H^p(X,\mathcal F)$ une structure de $A$-module, et l'opérateur $\partial$
est $A$-linéaire. En outre, pour deux entiers positifs quelconques $p$, $q$,
et deux $\mathcal O_X$-Modules $\mathcal F$, $\mathcal G$, on a un
homomorphisme de $A$-modules, dit \emph{cup-produit}
\[
\label{0.12.1.2.1-fr}
\tag{12.1.2.1}
H^p(X,\mathcal F)\otimes_A H^q(X,\mathcal G)
\longrightarrow H^{p+q}(X,\mathcal F\otimes_{\mathcal O_X}\mathcal G)
\]
(G, II, 6.6). Ces homomorphismes font de la somme directe $S$ des
$H^p(X,\mathcal O_X)$ (pour $p\geq0$) une $A$-algèbre graduée
anticommutative et de la somme directe des $H^p(X,\mathcal F)$ un $S$-module
gradué.

Pour tout \emph{recouvrement ouvert} $\mathfrak U$ de $X$, nous désignerons
toujours par $\check C^\bullet(\mathfrak U,\mathcal F)$ (contrairement à
(G, II, 5.1)) le complexe des cochaînes \emph{alternées} du nerf de
$\mathfrak U$ à valeurs dans le système de coefficients
$\Gamma(U_\alpha,\mathcal F)$. Il est clair que
$\check C^\bullet(\mathfrak U,\mathcal F)$ est un $A$-module gradué, donc les
groupes de cohomologie $\check H^p(\mathfrak U,\mathcal F)$ de ce complexe
sont munis d'une structure de $A$-module; en outre, les applications
canoniques
\[
\check H^p(\mathfrak U,\mathcal F)\longrightarrow H^p(X,\mathcal F)
\]
(G, II, 5.4) sont $A$-linéaires.
\end{env}

\begin{env}[12.1.3]
\label{0.12.1.3-fr}
Soit $(X',\mathcal O_{X'})$ un second espace annelé, et soit
$f=(\psi,\theta)$ un morphisme de $X'$ dans $X$.

Posons $A'=\Gamma(X',\mathcal O_{X'})$; le $\psi$-morphisme $\theta$ définit
canoniquement un homomorphisme d'anneaux $A\to A'$. Soient $\mathcal F$ un
$\mathcal O_X$-Module, $\mathcal F'$ un $\mathcal O_{X'}$-Module; pour tout
$f$-morphisme $u:\mathcal F\to\mathcal F'$
(\hyperref[0.4.4.1-fr]{0\textsubscript{I}, 4.4.1}), nous allons voir qu'on
peut définir pour tout $p\geq0$, un \emph{di-homomorphisme}
\[
\label{0.12.1.3.1-fr}
\tag{12.1.3.1}
u_p:H^p(X,\mathcal F)\longrightarrow H^p(X',\mathcal F').
\]
En effet, comme $\psi^*$ est exact dans la catégorie des faisceaux de groupes
abéliens sur $X$, $\mathcal F\to H^\bullet(X',\psi^*(\mathcal F))$ est un
$\delta$-foncteur dans cette catégorie, et on sait qu'on a un homomorphisme
canonique de $\delta$-foncteurs
\[
\label{0.12.1.3.2-fr}
\tag{12.1.3.2}
H^\bullet(X,\mathcal F)\longrightarrow H^\bullet(X',\psi^*(\mathcal F))
\]
uniquement déterminé par la condition de se réduire à l'homomorphisme
canonique $\Gamma(X,\mathcal F)\to\Gamma(X',\psi^*(\mathcal F))$ en degré
0 (T, 3.2.2). En outre, tout élément de $A$ détermine un endomorphisme $\mu$
de $\Gamma(X,\mathcal F)$ et un endomorphisme $\mu'$ de
$\Gamma(X',\psi^*(\mathcal F))$ tels que le diagramme
\[
\label{0.12.1.3.3-fr}
\tag{12.1.3.3}
\xymatrix{
\Gamma(X,\mathcal F)\ar[r]\ar[d]_\mu
&\Gamma(X',\psi^*(\mathcal F))\ar[d]^{\mu'}\\
\Gamma(X,\mathcal F)\ar[r]
&\Gamma(X',\psi^*(\mathcal F))
}
\]
soit commutatif; par la propriété d'unicité de prolongement des morphismes
pour les foncteurs cohomologiques universels (T, 2.2), on en déduit des
prolongements uniques de $\mu$ et $\mu'$ à la cohomologie rendant commutatifs
les diagrammes analogues à (\hyperref[0.12.1.3.3-fr]{12.1.3.3}), ce qui
signifie que (\hyperref[0.12.1.3.2-fr]{12.1.3.2}) est un homomorphisme de
$A$-modules. Notons maintenant que l'on a
\[
f^*(\mathcal F)=\psi^*(\mathcal F)
\otimes_{\psi^*(\mathcal O_X)}\mathcal O_{X'},
\]
et que l'on a donc un di-homomorphisme canonique
$\psi^*(\mathcal F)\to f^*(\mathcal F)$ du
$\psi^*(\mathcal O_X)$-Module $\psi^*(\mathcal F)$ dans le
$\mathcal O_{X'}$-Module $f^*(\mathcal F)$. Par fonctorialité, on en déduit
donc un di-homomorphisme fonctoriel
\[
\label{0.12.1.3.4-fr}
\tag{12.1.3.4}
H^p(X',\psi^*(\mathcal F))\longrightarrow H^p(X',f^*(\mathcal F))
\]
les anneaux correspondants étant $A$ et $A'$; en composant ce
di-homomorphisme avec (\hyperref[0.12.1.3.2-fr]{12.1.3.2}), on obtient un
di-homomorphisme canonique fonctoriel en $\mathcal F$
\[
\label{0.12.1.3.5-fr}
\tag{12.1.3.5}
\theta_p:H^p(X,\mathcal F)\longrightarrow H^p(X',f^*(\mathcal F)).
\]
Enfin, par fonctorialité, on déduit de l'homomorphisme
$u^\sharp:f^*(\mathcal F)\to\mathcal F'$ un homomorphisme de $A'$-modules
$H^p(X',f^*(\mathcal F))\to H^p(X',\mathcal F')$, qui, composé avec
(\hyperref[0.12.1.3.5-fr]{12.1.3.5}), donne
(\hyperref[0.12.1.3.1-fr]{12.1.3.1}).

Soit $f'=(\psi',\theta'):X''\to X'$ un second morphisme d'espaces annelés,
$f''=ff'$ le morphisme composé. Compte tenu de la permutabilité du foncteur
$\psi'^*$ et du produit tensoriel
(\hyperref[0.4.3.3-fr]{0\textsubscript{I}, 4.3.3}), on vérifie aussitôt que le
composé du di-homomorphisme
$H^p(X',f^*(\mathcal F))\to H^p(X'',f'^*(f^*(\mathcal F)))$ et de
(\hyperref[0.12.1.3.5-fr]{12.1.3.5}) est le di-homomorphisme correspondant
$H^p(X,\mathcal F)\to H^p(X'',f''^*(\mathcal F))$.
\end{env}

\begin{env}[12.1.4]
\label{0.12.1.4-fr}
Une définition directe de l'homomorphisme
(\hyperref[0.12.1.3.2-fr]{12.1.3.2}) peut s'obtenir de la façon suivante : on
considère une résolution injective $\mathcal L^\bullet=(\mathcal L^i)$ de
$\mathcal F$ formée de faisceaux de groupes abéliens sur $X$; comme le
foncteur $\psi^*$ est exact, $\psi^*(\mathcal L^\bullet)$ est une résolution
de $\psi^*(\mathcal F)$ formée de faisceaux sur $X'$. Si
$\mathcal L'^\bullet=(\mathcal L'^i)$ est une résolution injective de
$\psi^*(\mathcal F)$ dans la catégorie des faisceaux de groupes abéliens sur
$X'$, il y a donc un morphisme
$\psi^*(\mathcal L^\bullet)\to\mathcal L'^\bullet$ de complexes de faisceaux
de groupes abéliens, compatible avec les augmentations (M, V, 1.1.a)), bien
déterminé à une homotopie près. On en déduit des homomorphismes
\[
\Gamma(X,\mathcal L^\bullet)\longrightarrow
\Gamma(X',\psi^*(\mathcal L^\bullet))\longrightarrow
\Gamma(X',\mathcal L'^\bullet)
\]
de complexes de groupes abéliens, dont le composé, par passage à la
cohomologie donne un morphisme de $\delta$-foncteurs
$H^\bullet(X,\mathcal F)\to H^\bullet(X',\psi^*(\mathcal F))$; comme il
coïncide avec (\hyperref[0.12.1.3.2-fr]{12.1.3.2}) en degré 0, il lui est
identique (T, 2.2).

Considérons maintenant un \emph{recouvrement ouvert}
$\mathfrak U=(U_\alpha)$ de $X$, et soit
$\mathfrak U'=(f^{-1}(U_\alpha))$ le recouvrement ouvert de $X'$, image
réciproque de $\mathfrak U$. Les homomorphismes canoniques
$\Gamma(V,\mathcal F)\to\Gamma(f^{-1}(V),f^*(\mathcal F))$ pour tout ouvert
$V$ de $X$ définissent aussitôt (cf. G, II, 5.1) un homomorphisme de complexes
$\check C^\bullet(\mathfrak U,\mathcal F)\to
\check C^\bullet(\mathfrak U',f^*(\mathcal F))$, d'où des homomorphismes
canoniques
\[
\label{0.12.1.4.1-fr}
\tag{12.1.4.1}
\theta_p:\check H^p(\mathfrak U,\mathcal F)
\longrightarrow\check H^p(\mathfrak U',f^*(\mathcal F)).
\]

